\subsection*{MT-M36. The literal section average is zero;
normalized compact projections and genuine idelic cutoffs}

This is a targeted measure and normalization check of
Masullo's formula (11.3.2.2), original PDF p.~495.
The formula uses ordinary indicator functions and the
ordinary Haar integral on the idelic fundamental domain.
It does not display a Dirac distribution or a rescaled
limit of neighborhoods.  Under precisely that
interpretation its value is identically zero.
We prove this, and then give two legitimate replacements
with their exact scopes and trace properties.

\paragraph{The source's measure conventions.}
Write $\mathbb I_F=\mathrm{GL}_1(\mathbf A_F)$,
$\mathbb I_F^1=\ker|\cdot|_F$, and let
$\tau:\mathbf R_{>0}\to\mathbb I_F$ be a continuous
homomorphic section of the norm.  The source defines
\[
 E_F^1=\bigcup_{i=1}^{h_F}
              g_i(E_0\times\widehat{\mathcal O}_F^\times)
                         \subset\mathbb I_F^1
 \tag{M36.1}
\]
in Definition A.2.3.6, (A.2.3.14), original p.~549,
and explicitly gives this set the subspace topology.
It is not assigned the discrete topology.  Definition
A.3.2.1, pp.~559--560, uses Haar measures on ideles
and their class groups; the counting measure is on
$F^\times$.  The same distinction is explicit in
the compact-character integration of \S2.2.5.
The normalization (A.3.2.1), p.~560, gives volume one
on $C_{\mathbf Q}^1$.  The general measure discussion
in \S12.2, p.~504, is ordinary completed Radon/Lebesgue
integration of measurable functions.

The displayed two-variable function (11.3.2.2) is
therefore, with its suppressed Haar differential written,
\[
 B_0(x,y)=\sum_{\xi\in F^\times}\int_{E_F^1}
    \mathbf1_{\tau(\mathbf R_{>0})}(\xi gx)
    \mathbf1_{\tau(\mathbf R_{>0})}(\xi gy)\,d\mu_1(g).
 \tag{M36.2}
\]
The tensor symbol in the source just separates the two
named variables of these indicator functions.
For every $x,y\in\mathbb I_F$,
\[
                              B_0(x,y)=0.
 \tag{M36.3}
\]

\begin{proof}
Haar measure on $\mathbb I_F^1$ has no atoms.
Indeed this group contains an infinite compact subgroup:
embed $\mathcal O_{F_v}^\times$ at any one finite place
and use $1$ at every other place.  If one singleton had
positive Haar measure $c$, every singleton would have
the same measure, and $N$ distinct points in this compact
subgroup would have measure $Nc$ for every $N$.
This contradicts finiteness of a Radon measure on a
compact set.  Consequently every singleton in $E_F^1$
is null as well.  Equivalently, the quotient Haar
measure pulled back to a fundamental domain is atomless;
the discrete quotient by $F^\times$ does not change this
fact.  For $F=\mathbf Q$ this is particularly explicit:
$E_{\mathbf Q}^1=\{1\}\times\widehat{\mathbf Z}^\times$,
with its usual compact, atomless, volume-one Haar measure.

The section image is closed, hence Borel: it is the
fixed-point set of the continuous map
$z\mapsto\tau(|z|_F)$ on the Hausdorff group
$\mathbb I_F$.  Fix $\xi$ and $x$.  If
$\xi gx=\tau(t)$ with $g\in E_F^1$, the product formula
and $|g|_F=1$ force $t=|x|_F$.  Hence the only possible
value of $g$ is
\[
                       g=\xi^{-1}\tau(|x|_F)x^{-1}.
 \tag{M36.4}
\]
The first indicator in the integrand of (M36.2) is
supported on at most this one point.  Its product with
the second indicator has integral zero.  The field
$F$ is countable, so (M36.2) is a countable sum of zeros.
This proves (M36.3), without any interchange of
conditionally convergent expressions.
\end{proof}

\paragraph{The precise normalization failure.}
The source's indentation condition (12.2.1.2), original
p.~505, requires an orbit integral equal to one.
For a discrete subgroup $H$ this reads
$\sum_{\xi\in H}b(\xi x)=1$, with $b\ge0$.
Its two-variable lift is the averaged product
\[
                    B_b(x,y)=\sum_{\xi\in H}b(\xi x)b(\xi y),
 \tag{M36.5}
\]
as in (12.2.2.1), p.~508.  Nonnegativity and reindexing
give the necessary orbit identity
\[
 \sum_{\eta\in H}B_b(x,\eta x)
    =\left(\sum_{\xi\in H}b(\xi x)\right)^2=1.
 \tag{M36.6}
\]
Thus the zero function in (M36.3) cannot provide this
normalized lift.  The source's actual lifting rule
(12.2.2.2), p.~508, is multiplication of the pulled-back
kernel by the two-variable factor.  Inserting (M36.2)
into that rule produces the zero kernel.

This also gives a concrete nonzero test for the literal
interface in Theorem 11.3.2.1.  Take $F=\mathbf Q$,
the trivial character, $S=\{\infty\}$, $s=2$, and any
allowed interval $I$.  Fix a rational prime $p$ and
choose a nonnegative $\psi\in C_c^\infty((0,\infty))$
supported in a small neighborhood of $\log p$, containing
no other element of $\{n\log p:n\in\mathbf Z\}$, with
$\psi(\log p)>0$.  The test
\[
                             \varphi(g)=\psi(\log|g|)
 \tag{M36.7}
\]
is compactly supported on $C_{\mathbf Q}$, is invariant
under every finite local unit group, and satisfies all
the displayed ramification and support hypotheses.
For a local uniformizer $\pi_p$, whose idele norm is
$p^{-1}$, the right side of (11.3.2.1) is
\[
       \frac{\log p}{\mu_p(\mathbf Z_p^\times)}
                           \psi(\log p)>0,
 \tag{M36.8}
\]
because only the power $\pi_p^{-1}$ contributes.
Under the stated multiplication rule the left side
with (M36.2) is zero.  This refutes the literal
normalization/interface, not a reinterpretation using
a separately specified distributional regularization.

\paragraph{A normalized finite compact-level replacement.}
Let $K$ be an infinite compact metrizable profinite
abelian group with probability Haar measure; in the
actual rational example $K=C_{\mathbf Q}^1\simeq
\widehat{\mathbf Z}^\times$.  For an open subgroup
$U\subset K$, put $m_U=\mu(U)=1/[K:U]$ and define
\[
 a_U(k)=\frac{\mathbf1_U(k)}{m_U},\qquad
 B_U(k,l)=\int_Ka_U(rk)a_U(rl)\,d\mu(r).
 \tag{M36.9}
\]
Then
\[
 B_U(k,l)=\frac{\mathbf1_U(kl^{-1})}{m_U}
         =\sum_{\chi\in\widehat{K/U}}\chi(k)\overline{\chi(l)}.
 \tag{M36.10}
\]
The first equality follows because the two cosets of
$U$ in the integrand are either disjoint or identical;
in the latter case their intersection has measure $m_U$.
The second is the elementary character orthogonality
identity on the finite abelian group $K/U$.

Consequently the integral operator with kernel $B_U$
is exactly the orthogonal projection
\[
      (E_Uf)(k)=\frac1{m_U}\int_U f(ku)\,d\mu(u)
                    \quad\hbox{on }L^2(K).
 \tag{M36.11}
\]
Its range is the space of $U$-invariant functions,
and
\[
 \operatorname{rank}E_U=\operatorname{Tr}E_U=m_U^{-1},
       \quad\|E_U\|=1,\quad
       \|B_U\|_{L^2(K\times K)}^2=m_U^{-1}.
 \tag{M36.12}
\]
All these assertions follow either from (M36.10) or
from the orthonormal basis of normalized coset indicators.

If the neighborhood indicators are left unnormalized,
their averaged product is instead
\[
 \int_K\mathbf1_U(rk)\mathbf1_U(rl)\,d\mu(r)
       =m_U\mathbf1_U(kl^{-1})=m_U^2 B_U(k,l).
 \tag{M36.13}
\]
Its operator norm is $m_U^2$ and its trace is $m_U$.
Along a nested sequence of open subgroups $U_j$ forming
a neighborhood base at $1$, one has $m_{U_j}\to0$,
because their intersection is the Haar-null singleton.
Thus these unnormalized operators tend to zero even
in trace norm.  The normalized $E_{U_j}$ instead
converge strongly to the identity.  To prove the latter,
locally constant functions are dense in $L^2(K)$:
continuous functions are dense, and uniform continuity
approximates each by functions constant on the cosets
of a sufficiently small open subgroup.  Every such
function is fixed by all sufficiently late $E_{U_j}$,
and the projections have norm one.  Density proves
strong convergence.  There is no operator-norm
convergence to the identity, since each range is finite
dimensional whereas $L^2(K)$ is infinite dimensional.

These projections are a genuine normalized regularization
of the compact section.  Their diverging diagonal values
and Hilbert--Schmidt norms in (M36.12) also show why they
cannot be replaced by an ordinary point-indicator kernel.

\paragraph{Trace compatibility on the actual Hilbert modules.}
The same averaging defines orthogonal projections $E_U$
for any strongly continuous unitary representation of
$K$ on a Hilbert space $\mathcal H$.  They converge
strongly to the identity as $U$ shrinks to $1$:
\[
 \|E_Uv-v\|
       \le\sup_{u\in U}\|R(u)v-v\|\longrightarrow0.
 \tag{M36.14}
\]
This includes the compact action on the weighted idelic
Hilbert spaces and their closed invariant quotients:
the norm weight depends only on the radial variable.
If $T$ is trace class, then
\[
 \|E_UT-T\|_1\to0,\qquad
 \|TE_U-T\|_1\to0,\qquad
             \operatorname{Tr}(E_UT)\to\operatorname{Tr}(T).
 \tag{M36.15}
\]
For a complete proof, write a singular-value expansion
$T=\sum_ns_n u_n\otimes v_n^*$ with
$\sum_ns_n<\infty$ and unit vectors.  The first
trace-norm difference is at most
$\sum_ns_n\|(E_U-1)u_n\|$, which tends to zero by
(M36.14) and domination by $2s_n$.  For the second
use self-adjointness of $E_U$ and the same argument
on $v_n$.  Continuity of the trace gives the last claim.

There is exact stabilization for a test already at
finite compact level.  If a convolution test $\varphi$
is invariant under an open subgroup $U_0\subset K$,
its convolution operator $A_\varphi$ on any of these
modules satisfies
\[
        E_U A_\varphi=A_\varphi E_U=A_\varphi
                          \qquad(U\subset U_0).
 \tag{M36.16}
\]
Indeed translation of the test by any $u\in U_0$
does not change its convolution operator; integration
over $U$ gives both identities, using commutativity.
Thus if $A_\varphi$ is trace class its trace is preserved
exactly by all these finite-level projections.  On the
arithmetic quotients that trace-class hypothesis is
available unconditionally in the wide strips of MT-M27
and, with its explicit Hecke-GRH premise, in the narrow
positive-width strips of MT-M28.  Formula (M36.16) does
not supply that analytic hypothesis by itself.
For general number fields the compact group can have
connected archimedean factors; the finite-rank assertions
(M36.10)--(M36.12) here are expressly for the profinite
group, not for an untruncated archimedean character family.

\paragraph{A genuine discrete-idelic indentation and its orbital allocation.}
There is a different, also explicit, correction for
the discrete $\mathbf Q^\times$-normalization in
(M36.5)--(M36.6).  Use the standard section and set
\[
 D=\tau(\mathbf R_{>0})\,\widehat{\mathbf Z}^\times
       \subset\mathbb I_{\mathbf Q},\qquad b=\mathbf1_D.
 \tag{M36.17}
\]
The set $D$ is an open subgroup, hence also closed.
Every idele has a unique expression $q d$ with
$q\in\mathbf Q^\times$ and $d\in D$: cancel its
finite valuations by a rational number and then fix
the sign of its real component.  Uniqueness follows
from $\mathbf Q^\times\cap D=\{1\}$.
The map $D\to C_{\mathbf Q}$ is a homeomorphism;
openness follows from the quotient map and the
openness of $D$.  Therefore $b$ is continuous and
nonnegative, satisfies $\sum_qb(qx)=1$, and its support
meets the inverse image of any compact subset of
$C_{\mathbf Q}$ in a compact set.  These are exactly
the three indentation conditions of Lemma 12.2.1.1.
No auxiliary connected-domain assumption is needed
for this direct verification.

For $x\in D$, its two-variable lift has
\[
                   B_b(x,qx)=\mathbf1_{q=1}.
 \tag{M36.18}
\]
Thus it provides the normalized total orbital integral,
but places all of that diagonal mass in the identity
translate.  This allocation is not forced by the
normalization.  For a rational prime $p$ and
$0\le t\le1$, another valid indentation is
\[
                    b_t=(1-t)\mathbf1_D+t\mathbf1_{pD}.
 \tag{M36.19}
\]
Its support is a union of two clopen translates of
$D$, it has the same compactness property, and
$\sum_qb_t(qx)=1$.  Direct substitution, for $x\in D$,
gives
\[
 B_{b_t}(x,qx)=
 \begin{cases}
       (1-t)^2+t^2,&q=1,\\
       t(1-t),&q=p\text{ or }p^{-1},\\
       0,&\text{otherwise}.
 \end{cases}
 \tag{M36.20}
\]
The sum of these values is one.  Consequently, when
an integrable quotient diagonal is pulled back and
multiplied by $B_{b_t}$, its total orbital integral is
unchanged, while its individual rational-translate
contributions vary with $t$.  This follows directly
from (M36.20), since $x$ and $qx$ represent the same
class; no trace-limit interchange is used.

The normalized compact projection and the genuine
discrete-idelic indentation solve different problems.
Neither calculation identifies the source's desired
prime-by-prime geometric mass with the local Poisson
trace of MT-M32.  A correct comparison must specify
its normalized lift and prove the required orbital
allocation and limiting identity.  Positivity or the
conditional Hecke-GRH endpoint does not follow from
the measure repair alone.

\paragraph{Audit scope.}
Original PDF pp.~495, 505, 508, 549 and 560 were
rendered and visually read.  The definitions and
normalization text on pp.~504--506 and 559--560 were
also read from the supplied mathematical text.
This is not an audit of all the kernel theorems in
Chapter 12.  The literal failure (M36.3) and all the
replacement constructions above are unconditional.
The endpoint $M_H$ remains conditional throughout.

For cote 1, scans 79--80 and 111--112, the distinction
is between a normalized lift of a quotient trace and
its claimed local closed-point contributions.  The
former is available in the two explicit forms above;
the latter still requires a canonical allocation and
a proved geometric comparison.
